% mage-auto-heading: the-capability
% mage-auto-heading: when-this-stack-earns-its-keep
% mage-auto-heading: the-composition
% mage-auto-heading: constituent-moves
% mage-auto-heading: why-these-travel-together
% Authoritative LaTeX manuscript; Markdown is a generated compatibility projection.

\section{The Capability}\label{heading:the-capability}

\textbf{State the engineering obligations, identify the population they apply to, and produce evidence strong enough for each property.}

\section{When This Stack Earns Its Keep}\label{heading:when-this-stack-earns-its-keep}

Reach for it when the important question shifts from

% <!-- inline-quote -->

\begin{quote}
Did the tests pass?
\end{quote}

to

% <!-- inline-quote -->

\begin{quote}
Have we named the relevant obligations, discharged each one appropriately, and shown the population is covered?
\end{quote}

That shift arrives whenever a property has to hold across a whole class of things --- every mutator, every config field, every seam --- and a green suite on a handful of examples no longer settles it.

\section{The Composition}\label{heading:the-composition}

\begin{figure}[htbp]
\centering
\includegraphics[width=\linewidth]{book/latex/assets/assurance-stack-0fad62ad9e.pdf}
\caption{The assurance composition. A SPEC states the obligation; a CENSUS establishes the population it applies to; discharge fans out to the evidence each obligation deserves --- TEST for examples, LINT for structure, PROVE for semantics --- and all three lanes converge on COVERAGE, which joins each obligation back to its evidence so omissions show. Solid path: the required spine.}
\label{assurance-stack}
\end{figure}

\section{Constituent Moves}\label{heading:constituent-moves}

{\footnotesize
{\def\LTcaptype{none} % do not increment counter
\begin{longtable}[]{@{}
  >{\raggedright\arraybackslash}p{(\linewidth - 2\tabcolsep) * \real{0.5000}}
  >{\raggedright\arraybackslash}p{(\linewidth - 2\tabcolsep) * \real{0.5000}}@{}}
\toprule\noalign{}
\begin{minipage}[b]{\linewidth}\raggedright
Move
\end{minipage} & \begin{minipage}[b]{\linewidth}\raggedright
Role
\end{minipage} \\
\midrule\noalign{}
\endhead
\bottomrule\noalign{}
\endlastfoot
\textbf{SPEC} & State the obligation. \\
\textbf{CENSUS} & Establish the population to which it applies. \\
\textbf{DISCHARGE} & Apply tests, lints, proofs, or other evidence appropriate to the property. \\
\textbf{COVER} & Join obligations back to evidence so omissions become visible. \\
\end{longtable}
}
}

\section{Why These Travel Together}\label{heading:why-these-travel-together}

A checker establishes a property only for what it checks. A proof says nothing about obligations left out of its population; a census with no evidence attached merely enumerates debt. Assurance therefore requires both evidence for each obligation and confidence that the relevant population is covered. Use evidence suited to the property---an example test where an example suffices, a structural lint where structure is decisive, a bounded proof where semantics demand it---and map that evidence back to the full population so omissions remain visible.

One rule runs throughout: check an obligation at the semantic level where the property becomes legible. A property about a document\textquotesingle s structure belongs in a check that reads structure, not one that scans bytes for it. This determines where the check belongs; it is not another mechanism in the stack.

\textbf{Mechanisms:} census-derived obligations · semantic validator · bounded proof · coverage gate
